\chapter{Conclusion}
\label{ch:Conclusion}
This researched analyzed \textit{Oblivious Random Access Memories (ORAMs)} with arbitrary element size.
It introduced ORAM including its general concepts, technical foundations and goals.
\textit{DUORAM}, as introduced by Vadapalli et al. \cite{vadapalli2023duoram}, was taken as an example.

Using the implementation from Vadapalli et al. as foundation, this research developed and implemented two approaches on how to scale DUORAM to arbitrary element sizes.
\textit{new\_input\_type} executes multiple reads or updates and aggregates its result to a single entry.
This is done by executing the already implemented operations multiple times.
\textit{gmp\_mpz\_class} uses the \textit{GNU Multi Precision Arithmetic Library (GMP)} to scale not just the entries but also the RDPFs.
This enables the user to store and retrieve elements of arbitrary size into a single entry.
Both approaches require the user to commit to a fixed entry size before compilation.

Both approaches are tested on 64, 128, and 256 bit.
Larger bit sizes cannot be executed due to technical restrictions of the author's machine.

\textit{new\_input\_type} proved to be faster on average for both the preprocessing and the online phase.
Its communication methods are faster, and bit operations can be executed on sizes supported directly by the system.
Also, no serialization and deserialization methods need to be executed before communicating.
\textit{new\_input\_type} can utilize multithreading to perform reads and updates in parallel.
This pushes its performance for 128 and 256 bits to almost the performance of the base implementation by Vadapalli et al.
Performance, however, heavily deteriorates for low bandwidths.
Also, it has heavier storage requirements for depths greater 15 than \textit{gmp\_mpz\_class} and a higher consumption of RDPF-triples.

\textit{gmp\_mpz\_class} is faster for lower bandwidths as messages are sent in smaller packages.
As consequently more messages need to be sent, performance deteriorates for higher latencies.
The underlying datatype \code{mpz\_class} has a memory overhead of factor $1.5$ on average for trivial datatypes of the same size.
It has, however, lower storage requirements than \textit{new\_input\_type} as fewer RDPF-triples need to be generated due to lower consumption during the online phase.
Also, bit operations on \code{mpz\_class} values are slower than bit operations on trivial datatypes of the same size.
This was proven true by testing for XOR and bit shifting operations, which are heavily used throughout the implementation.
Finally, communication is slower due to the need of serialization and deserialization while sending or receiving.

Therefore, \textit{new\_input\_type} is preferable on machines with larger storage and environments of high bandwidth and varying latency.
\textit{gmp\_mpz\_class} proved to be the better approach for machines with smaller storage, varying bandwidth, and smaller latency.

GMP in general provides useful datatypes and classes for this research's goal of scaling \textit{DUORAM}.
\code{mpz\_t} represents an integer of arbitrary size.
Yet, it does not provide automatic memory management which makes it uneasy to handle in a complex program without compromising type safety.
\code{mpz\_class} holds a single attribute of type \code{mpz\_t} and provides automatic memory management.
It has, as previously shown, drawbacks in its storage overhead.
Also, bit operations on \code{mpz\_t} and \code{mpz\_class} are slower than on trivial datatypes of the same size.
Another practical drawback is its characteristics during compilation.
Since both \code{mpz\_t} and \code{mpz\_class} are evaluated during runtime they cannot be set as template parameters which needs the properties' arguments during preprocessing.
However, as \code{mpz\_t} and \code{mpz\_class} provide arbitrarily scalable integers, bit operations of acceptable performance and a usable documentation they proved the most useful for this research.

The author recommends further research.
First, one can test the existing implementation on more performant machines.
This includes testing for larger bit sizes, larger numbers of preprocessed RDPF-triples, RDPF-triples of greater depth, and more items retrieved or updated during the online phase.
Second, the existing approaches can be tested for different datatypes.
Possible alternatives are \textit{Boost.Multiprecision} or \textit{Fast Library for Number Theory (FLINT)}.
Finally, other approaches can be analyzed that could not be implemented or described for time reasons.
One approach of this type could use two DCFs (\ref{cha:DistributedComparisonFunctions}) per operation.
These DCFs define an interval in the database in which one reads or updates her entries.